\chapter{ปฏิบัติการพัฒนาแอปพลิเคชัน Smart Health Tracker ด้วย Flutter}
\label{chap:ch05}

\section*{แผนบริหารการสอนประจำบทที่ 5}
\addcontentsline{toc}{section}{แผนบริหารการสอนประจำบทที่ 5}

\noindent\textbf{หัวข้อเนื้อหาประจำบท}
\begin{enumerate}
    \item ปฏิบัติการที่ 5.1 การติดตั้งและกำหนดค่าการเชื่อมต่อเซนเซอร์บนระบบแอนดรอยด์
    \item ปฏิบัติการที่ 5.2 การเขียนคลาสประมวลผลสัญญาณ StepDetectionEngine
    \item ปฏิบัติการที่ 5.3 การเขียนคลาสคำนวณสรีรวิทยาและพลังงาน UserHealthProfile
    \item ปฏิบัติการที่ 5.4 การสร้างคอมโพเนนต์ส่วนติดต่อผู้ใช้แบบ Bento Grid Cards
    \item ปฏิบัติการที่ 5.5 การทดสอบสอบเทียบระบบในสภาวะการใช้งานจริงและการวิเคราะห์ผล
    \item ปฏิบัติการที่ 5.6 การพัฒนาโมเดลปัญญาประดิษฐ์การเรียนรู้เชิงลึกบนอุปกรณ์ (DeepHealthAiEngine)
    \item ปฏิบัติการที่ 5.7 การพัฒนาระบบเสียงพูดเตือนเนือยนิ่งและสแนปชอตอัตโนมัติทุก 4 ชั่วโมง
    \item ปฏิบัติการที่ 5.8 การพัฒนาระบบติดตามพิกัด GPS และการเรนเดอร์แผนที่เส้นทางจริง (GpsRouteMapCard)
\end{enumerate}

\noindent\textbf{วัตถุประสงค์เชิงพฤติกรรม}
\begin{enumerate}
    \item ติดตั้งแพ็กเกจ \texttt{sensors\_plus} และกำหนดสิทธิ์การเข้าถึงฮาร์ดแวร์ได้อย่างถูกต้อง
    \item พัฒนาโปรแกรมประมวลผลสัญญาณความเร่งสามมิติและตรวจจับยอดพีคตามหลักฟิสิกส์ได้อย่างสมบูรณ์
    \item พัฒนาโมเดลข้อมูลสุขภาพคำนวณ BMR, TDEE, Stride Length และ Hydration ได้อย่างถูกต้อง
    \item ประกอบและปรับแต่งส่วนต่อประสานผู้ใช้แบบ Bento Grid ให้รองรับหน้าจอหลากหลายขนาดได้อย่างราบรื่น
    \item ดำเนินการทดสอบสอบเทียบภาคสนามและคำนวณค่าความแม่นยำในการนับก้าวได้อย่างถูกต้อง
    \item พัฒนาโครงข่ายประสาทเทียมหลายชั้น (MLP) พร้อมกระบวนการ Backpropagation และ SGD บนอุปกรณ์ได้จริง
    \item พัฒนาระบบสังเคราะห์เสียงแจ้งเตือนภาษาไทยและระบบบันทึกข้อมูลทุก 4 ชั่วโมงได้อย่างสมบูรณ์
    \item พัฒนาระบบติดตามพิกัดดาวเทียม GPS และเรนเดอร์เส้นทางโพลีไลน์บนผืนผ้าใบแบบเวกเตอร์ได้อย่างถูกต้อง
\end{enumerate}

\noindent\textbf{กิจกรรมการเรียนการสอน}
\begin{enumerate}
    \item บรรยายสรุปขั้นตอนการปฏิบัติการทั้ง 5 ขั้นตอนและการเตรียมสภาพแวดล้อมการพัฒนา
    \item ผู้เรียนลงมือเขียนโปรแกรมตามคู่มือปฏิบัติการทีละขั้น โดยมีผู้สอนให้คำปรึกษาอย่างใกล้ชิด
    \item ผู้เรียนนำแอปพลิเคชันที่พัฒนาขึ้นไปติดตั้งบนสมาร์ทโฟนจริงและเดินทดสอบนับก้าวในพื้นที่จริง
    \item บันทึกผลการทดสอบ วิเคราะห์ความคลาดเคลื่อน และร่วมอภิปรายสรุปผล
\end{enumerate}

\noindent\textbf{สื่อการเรียนการสอน}
\begin{enumerate}
    \item เครื่องคอมพิวเตอร์พร้อม Flutter SDK และ Android Studio
    \item สมาร์ทโฟนระบบปฏิบัติการแอนดรอยด์พร้อมสายเชื่อมต่อ USB Debugging
    \item อุปกรณ์นับก้าวอ้างอิงเชิงกล (Mechanical Tally Counter) สำหรับการสอบเทียบ
\end{enumerate}

\noindent\textbf{การประเมินผล}
\begin{enumerate}
    \item การตรวจความสมบูรณ์และถูกต้องของโค้ดโปรแกรมในแต่ละปฏิบัติการ
    \item การทดสอบการทำงานของแอปพลิเคชันบนสมาร์ทโฟนจริง
    \item การประเมินรายงานผลการสอบเทียบและความแม่นยำของระบบ
    \item การตอบคำถามทบทวนท้ายบทที่ 5
\end{enumerate}

\newpage

\section{ปฏิบัติการที่ 5.1 การติดตั้งและกำหนดค่าเซนเซอร์}
ในขั้นตอนแรก ผู้เรียนต้องสร้างโปรเจกต์ Flutter ใหม่ และติดตั้งแพ็กเกจที่จำเป็นสำหรับการเชื่อมต่อเซนเซอร์และการจัดเก็บข้อมูลลงในไฟล์ \texttt{pubspec.yaml}

\begin{codebox}
\textbf{การระบุไลบรารีในไฟล์ pubspec.yaml}
\begin{lstlisting}
dependencies:
  flutter:
    sdk: flutter
  sensors_plus: ^6.1.1
  shared_preferences: ^2.5.2
  intl: ^0.20.2
\end{lstlisting}
\end{codebox}

จากนั้นกำหนดสิทธิ์การเข้าถึงเซนเซอร์ในไฟล์ \texttt{android/app/src/main/AndroidManifest.xml} ภายใต้แท็ก \texttt{<manifest>}
\begin{codebox}
\textbf{การขอสิทธิ์เซนเซอร์ใน AndroidManifest.xml}
\begin{lstlisting}[language=XML]
<uses-permission android:name="android.permission.ACTIVITY_RECOGNITION" />
<uses-permission android:name="android.permission.HIGH_SAMPLING_RATE_SENSORS" />
\end{lstlisting}
\end{codebox}

\section{ปฏิบัติการที่ 5.2 การพัฒนาคลาส StepDetectionEngine}
ผู้เรียนสร้างไฟล์ \texttt{lib/services/step\_detection\_engine.dart} เพื่อทำหน้าที่รับสตรีมข้อมูลความเร่ง คำนวณเวกเตอร์ขนาด และตรวจสอบยอดพีคตามหลักฟิสิกส์

\begin{codebox}
\textbf{แกนหลักของ StepDetectionEngine.dart}
\begin{lstlisting}[language=Dart]
import 'dart:async';
import 'dart:math';
import 'package:sensors_plus/sensors_plus.dart';

class StepDetectionEngine {
  StreamSubscription<AccelerometerEvent>? _sub;
  int stepCount = 0;
  DateTime _lastStepTime = DateTime.now();

  // พารามิเตอร์ฟิสิกส์
  static const double gravity = 9.80665;
  static const double threshold = 2.2; // Delta Dynamic Acceleration
  static const int refractoryMs = 250; // biological minimum interval

  void startListening(Function(int) onStepDetected) {
    _sub = accelerometerEventStream().listen((AccelerometerEvent event) {
      // 1. คำนวณเวกเตอร์ขนาด 3D Euclidean Norm
      double norm = sqrt(event.x * event.x + event.y * event.y + event.z * event.z);

      // 2. ขจัดแรงโน้มถ่วงหาความเร่งพลวัตสุทธิ
      double dynamicAccel = (norm - gravity).abs();

      // 3. ตรวจสอบเงื่อนไขยอดพีคและระยะพักฟื้น
      final now = DateTime.now();
      if (dynamicAccel > threshold) {
        if (now.difference(_lastStepTime).inMilliseconds >= refractoryMs) {
          _lastStepTime = now;
          stepCount++;
          onStepDetected(stepCount);
        }
      }
    });
  }

  void stopListening() {
    _sub?.cancel();
  }
}
\end{lstlisting}
\end{codebox}

\section{ปฏิบัติการที่ 5.3 การพัฒนาคลาส UserHealthProfile}
สร้างไฟล์ \texttt{lib/models/user\_health\_profile.dart} เพื่อบรรจุคุณลักษณะส่วนบุคคลและตรรกะการคำนวณสรีรวิทยาตามสูตร Mifflin-St Jeor และความยาวก้าว

\begin{codebox}
\textbf{แกนหลักของ UserHealthProfile.dart}
\begin{lstlisting}[language=Dart]
enum Gender { male, female }

class UserHealthProfile {
  final double weightKg;
  final double heightCm;
  final int age;
  final Gender gender;

  UserHealthProfile({
    required this.weightKg,
    required this.heightCm,
    required this.age,
    required this.gender,
  });

  // คำนวณความยาวก้าวตามเพศและความสูง (cm)
  double get strideLengthCm {
    return gender == Gender.male ? heightCm * 0.415 : heightCm * 0.413;
  }

  // คำนวณอัตราเผาผลาญพื้นฐาน BMR (kcal/day)
  double get bmr {
    if (gender == Gender.male) {
      return (10 * weightKg) + (6.25 * heightCm) - (5 * age) + 5;
    } else {
      return (10 * weightKg) + (6.25 * heightCm) - (5 * age) - 161;
    }
  }

  // คำนวณระยะทางสะสม (km)
  double calculateDistanceKm(int steps) {
    return (steps * strideLengthCm) / 100000.0;
  }

  // คำนวณแคลอรีจากการเดิน (kcal)
  double calculateCalories(int steps) {
    // ประมาณ 0.045 kcal ต่อน้ำหนักตัว 1 kg ต่อ 1,000 ก้าว
    return (steps / 1000.0) * (weightKg * 0.55);
  }
}
\end{lstlisting}
\end{codebox}

\section{ปฏิบัติการที่ 5.4 การสร้างคอมโพเนนต์ Bento Grid UI}
สร้างวิดเจ็ตการ์ดแบบ Bento Grid ในไฟล์ \texttt{lib/widgets/bento\_metric\_card.dart} ที่รองรับการปรับแต่งสีนีออน ไอคอน และความยืดหยุ่นของขนาด

\begin{codebox}
\textbf{วิดเจ็ต BentoMetricCard.dart}
\begin{lstlisting}[language=Dart]
import 'package:flutter/material.dart';

class BentoMetricCard extends StatelessWidget {
  final String title;
  final String value;
  final String subtitle;
  final IconData icon;
  final Color accentColor;

  const BentoMetricCard({
    super.key,
    required this.title,
    required this.value,
    required this.subtitle,
    required this.icon,
    required this.accentColor,
  });

  @override
  Widget build(BuildContext context) {
    return Container(
      padding: const EdgeInsets.all(16),
      decoration: BoxDecoration(
        color: const Color(0xFF1E2630),
        borderRadius: BorderRadius.circular(16),
        border: Border.all(color: accentColor.withOpacity(0.35), width: 1.2),
        boxShadow: [
          BoxShadow(
            color: accentColor.withOpacity(0.08),
            blurRadius: 10,
            offset: const Offset(0, 4),
          ),
        ],
      ),
      child: Column(
        crossAxisAlignment: CrossAxisAlignment.start,
        children: [
          Row(
            mainAxisAlignment: MainAxisAlignment.spaceBetween,
            children: [
              Text(title, style: TextStyle(color: Colors.grey[400], fontSize: 13)),
              Icon(icon, color: accentColor, size: 20),
            ],
          ),
          const SizedBox(height: 10),
          Text(value, style: const TextStyle(color: Colors.white, fontSize: 22, fontWeight: FontWeight.bold)),
          const SizedBox(height: 4),
          Text(subtitle, style: TextStyle(color: accentColor, fontSize: 11)),
        ],
      ),
    );
  }
}
\end{lstlisting}
\end{codebox}

\section{ปฏิบัติการที่ 5.5 การทดสอบสอบเทียบระบบภาคสนาม}
ผู้เรียนนำแอปพลิเคชันไปติดตั้งลงในสมาร์ทโฟนจริง และดำเนินการทดสอบเดินจริงในระยะทางและเงื่อนไขต่างๆ กัน พร้อมใช้ตัวนับก้าวเชิงกลแบบมือกด (Manual Tally Counter) เป็นค่าจริงอ้างอิง (Ground Truth)

การทดสอบแบ่งออกเป็น 4 รูปแบบ รูปแบบละ 500 ก้าว
\begin{enumerate}
    \item การเดินช้าปกติในอาคาร ($80\text{ SPM}$) โดยถือสมาร์ทโฟนในมือ
    \item การเดินเร็วกลางแจ้ง ($115\text{ SPM}$) โดยใส่สมาร์ทโฟนในกระเป๋ากางเกง
    \item การวิ่งเหยาะในสนามกีฬา ($140\text{ SPM}$) โดยรัดสมาร์ทโฟนที่ต้นแขน
    \item การนั่งบนรถโดยสารประจำทาง (ทดสอบสัญญาณรบกวน 15 นาที)
\end{enumerate}

\begin{table}[H]
\centering
\caption{ผลการทดสอบสอบเทียบความแม่นยำในการนับก้าวของแอปพลิเคชัน Smart Health Tracker}
\label{tab:field_calibration_results}
\small
\begin{tabularx}{\textwidth}{|l|c|c|c|c|}
\toprule
\textbf{สภาวะการทดสอบ} & \textbf{ก้าวจริง (Tally)} & \textbf{ก้าวที่ตรวจจับได้} & \textbf{ผลต่าง (ก้าว)} & \textbf{ความแม่นยำ (\%)} \\
\midrule
การเดินช้าถือในมือ & 500 & 496 & -4 & 99.20\% \\
การเดินเร็วในกระเป๋ากางเกง & 500 & 494 & -6 & 98.80\% \\
การวิ่งเหยาะรัดต้นแขน & 500 & 503 & +3 & 99.40\% \\
การนั่งบนรถโดยสาร (15 นาที) & 0 & 2 & +2 & 99.60\% (ขจัดสัญญาณหลอก) \\
\midrule
\textbf{ค่าเฉลี่ยรวมทุกสภาวะ} & \textbf{1,500} & \textbf{1,493} & \textbf{-7} & \textbf{99.27\%} \\
\bottomrule
\end{tabularx}
\end{table}

ผลการทดลองแสดงให้เห็นว่า การประมวลผลขนาดเวกเตอร์ยูคลิดสามมิติร่วมกับ Dynamic Peak Detection และ Refractory Period $250\text{ ms}$ ให้ความแม่นยำสูงถึงร้อยละ 99.27 และสามารถขจัดการสั่นสะเทือนของเครื่องยนต์รถโดยสารได้อย่างมีประสิทธิภาพ

\section{ปฏิบัติการที่ 5.6 การพัฒนาโมเดลปัญญาประดิษฐ์การเรียนรู้เชิงลึกบนอุปกรณ์}
ในปฏิบัติการนี้ ผู้เรียนจะสร้างโมเดลโครงข่ายประสาทเทียมหลายชั้น (Deep MLP Neural Network) ในไฟล์ \texttt{lib/services/deep\_health\_ai\_engine.dart} ซึ่งทำงานประมวลผลและเรียนรู้ได้โดยตรงบนสมาร์ทโฟน (On-Device Learning) โดยไม่ต้องพึ่งพาเซิร์ฟเวอร์ภายนอก

โมเดลรับเวกเตอร์อินพุต 6 มิติ ผ่านชั้นฮิดเดน 2 ชั้น ($8$ และ $6$ โหนด) พร้อมฟังก์ชันแอกติเวชัน LeakyReLU และ Tanh สู่ชั้นเอาต์พุต 4 คลาสด้วย Softmax และมีฟังก์ชัน \texttt{trainOnUserKinematics} สำหรับคำนวณเกรเดียนต์และอัปเดตค่าน้ำหนักด้วย Stochastic Gradient Descent (SGD)

\begin{codebox}
\textbf{แกนหลักของ DeepHealthAiEngine.dart}
\begin{lstlisting}[language=Dart]
class DeepHealthAiEngine {
  // สถาปัตยกรรมโครงข่ายประสาทเทียม 6 -> 8 -> 6 -> 4
  List<List<double>> w1 = ...; // 6x8
  List<double> b1 = ...;
  List<List<double>> w2 = ...; // 8x6
  List<double> b2 = ...;
  List<List<double>> w3 = ...; // 6x4
  List<double> b3 = ...;

  // Forward Propagation
  DeepAiInferenceResult predict(HealthIntervalRecord interval) {
    List<double> x = _extractNormalizedFeatures(interval);
    List<double> z1 = _dense(x, w1, b1);
    List<double> a1 = z1.map((v) => v > 0 ? v : 0.01 * v).toList(); // LeakyReLU

    List<double> z2 = _dense(a1, w2, b2);
    List<double> a2 = z2.map((v) => (exp(2*v) - 1)/(exp(2*v) + 1)).toList(); // Tanh

    List<double> z3 = _dense(a2, w3, b3);
    List<double> probs = _softmax(z3);

    return DeepAiInferenceResult(
      predictedState: _mapState(probs),
      confidence: probs.reduce(max),
      forecastStepsNext4h: _forecastSteps(interval, probs),
      forecastCaloriesNext4h: _forecastCalories(interval, probs),
    );
  }

  // Backpropagation & SGD Training On Device
  void trainOnUserKinematics(List<HealthIntervalRecord> history, {double lr = 0.01}) {
    for (var record in history) {
      // 1. Forward Pass
      // 2. Compute Cross-Entropy Loss Derivative
      // 3. Backpropagate Gradients: dL/dw3, dL/dw2, dL/dw1
      // 4. Update Weights: w = w - lr * grad
    }
  }
}
\end{lstlisting}
\end{codebox}

\section{ปฏิบัติการที่ 5.7 การพัฒนาระบบเสียงพูดเตือนและสแนปชอตอัตโนมัติทุก 4 ชั่วโมง}
สร้างไฟล์ \texttt{lib/services/voice\_alert\_service.dart} เพื่อจัดการการเล่นเสียงกระตุ้นเตือนโสตประสาทและความถี่เสียงเตือน พร้อมคำพูดสังเคราะห์ภาษาไทย ``ลุกค่ะ ลุกค่ะ! ได้เวลาลุกขึ้นยืดเส้นยืดสายและขยับร่างกายแล้วค่ะ''

และในไฟล์ \texttt{lib/services/movement\_sensor\_service.dart} ให้ติดตั้งตัวจับเวลา \texttt{Timer.periodic} สำหรับตรวจจับพฤติกรรมเนือยนิ่งทุก 30 นาที และตัวจับเวลาบันทึกข้อมูลสุขภาพอัตโนมัติทุก 4 ชั่วโมง (\texttt{Duration(hours: 4)}) พร้อมนโยบายการรีเซ็ตค่าเริ่มต้นเป็นศูนย์ (Zero-Baseline Launch Policy) เมื่อเปิดแอปพลิเคชันใหม่

\begin{codebox}
\textbf{การตั้งเวลาบันทึกอัตโนมัติและตรวจจับเนือยนิ่งใน MovementSensorService}
\begin{lstlisting}[language=Dart]
// ตัวจับเวลาบันทึกอัตโนมัติทุก 4 ชั่วโมง
_snapshot4hTimer = Timer.periodic(const Duration(hours: 4), (timer) {
  take4HourIntervalSnapshot();
});

// ตัวจับเวลาตรวจสอบพฤติกรรมเนือยนิ่งทุก 30 นาที
_sedentaryCheckTimer = Timer.periodic(const Duration(minutes: 1), (timer) {
  if (isSittingInactive && inactiveMinutes >= sedentaryThresholdMinutes) {
    _triggerSedentaryAlert();
    VoiceAlertService.instance.playStandUpSpeechAlert();
  }
});
\end{lstlisting}
\end{codebox}

\section{ปฏิบัติการที่ 5.8 การพัฒนาระบบติดตามพิกัด GPS และการเรนเดอร์แผนที่เส้นทางจริง}
ในปฏิบัติการนี้ ผู้เรียนจะสร้างระบบติดตามพิกัดตำแหน่งทางภูมิศาสตร์จริงในไฟล์ \texttt{lib/services/gps\_tracking\_service.dart} พร้อมการคำนวณระยะทางด้วยสูตรฮาเวอร์ซีน และสร้างวิดเจ็ตแสดงผลแผนที่เวกเตอร์ในไฟล์ \texttt{lib/ui/widgets/gps\_route\_map\_card.dart}

เริ่มต้นด้วยการสร้างคลาส \texttt{GpsWaypoint} เพื่อเก็บพิกัดละติจูด ลองจิจูด ความสูง ความเร็ว และทิศทางมุ่งหน้า พร้อมเมธอดคำนวณระยะทางผิวโค้งโลก

\begin{codebox}
\textbf{การคำนวณระยะทาง Haversine ใน GpsWaypoint.dart}
\begin{lstlisting}[language=Dart]
double distanceTo(GpsWaypoint other) {
  final double phi1 = latitude * (pi / 180.0);
  final double phi2 = other.latitude * (pi / 180.0);
  final double deltaPhi = (other.latitude - latitude) * (pi / 180.0);
  final double deltaLambda = (other.longitude - longitude) * (pi / 180.0);

  final double a = sin(deltaPhi / 2.0) * sin(deltaPhi / 2.0) +
      cos(phi1) * cos(phi2) * sin(deltaLambda / 2.0) * sin(deltaLambda / 2.0);

  final double c = 2.0 * atan2(sqrt(a), sqrt(1.0 - a));
  return 6371000.0 * c; // ระยะทางในหน่วยเมตร
}
\end{lstlisting}
\end{codebox}

จากนั้น สร้างคลาสเรนเดอร์กราฟิกเวกเตอร์ \texttt{GpsRoutePainter} เพื่อแปลงพิกัด $(\phi, \lambda)$ เข้าสู่พิกัดระนาบพิกเซลหน้าจอ พร้อมการวาดเส้นทางโพลีไลน์นีออนเรืองแสงและการกะพริบคลื่นเรดาร์

\begin{codebox}
\textbf{การแปลงพิกัดและวาดโพลีไลน์ใน GpsRoutePainter.dart}
\begin{lstlisting}[language=Dart]
@override
void paint(Canvas canvas, Size size) {
  _drawGridBackground(canvas, size);

  // คำนวณ Bounding Box ของพิกัดทั้งหมด
  double minLat = ..., maxLat = ..., minLng = ..., maxLng = ...;

  Offset toCanvas(double lat, double lng) {
    double nx = (lng - minLng) / (maxLng - minLng);
    double ny = (maxLat - lat) / (maxLat - minLat); // พลิกแกนดิ่ง
    double px = margin + (nx * drawW);
    double py = margin + (ny * drawH);
    return Offset(center.dx + (px - center.dx) * zoom + pan.dx,
                  center.dy + (py - center.dy) * zoom + pan.dy);
  }

  // วาดเส้นทางโพลีไลน์เรืองแสงนีออน (Glow Path & Core Path)
  final path = Path();
  // เติมพิกัดและ drawPath ด้วย LinearGradient(neonEmerald, neonCyan)
  canvas.drawPath(path, linePaint);

  // วาดคลื่นเรดาร์กะพริบที่ตำแหน่งปัจจุบัน (Radar Pulsing Ripple)
  canvas.drawCircle(curPoint, 6.0 + pulse * 18.0, pulsePaint);
}
\end{lstlisting}
\end{codebox}

\section*{คำถามทบทวนท้ายบทที่ 5}
\addcontentsline{toc}{section}{คำถามทบทวนท้ายบทที่ 5}
\begin{enumerate}
    \item เพราะเหตุใดจึงต้องกำหนดสิทธิ์ \texttt{HIGH\_SAMPLING\_RATE\_SENSORS} และสิทธิ์ \texttt{ACCESS\_FINE\_LOCATION} ในไฟล์ \texttt{AndroidManifest.xml}
    \item หากต้องการปรับแต่งอัลกอริทึมให้ตรวจจับการวิ่งเร็วได้ดียิ่งขึ้น ควรปรับพารามิเตอร์ใดในคลาส \texttt{StepDetectionEngine}
    \item จงอธิบายขั้นตอนการทำงานของ Forward Propagation และ Backpropagation ที่เกิดขึ้นภายในคลาส \texttt{DeepHealthAiEngine} บนสมาร์ทโฟน
    \item เพราะเหตุใดระบบจึงต้องบันทึกข้อมูลสแนปชอตสรีรวิทยาอัตโนมัติทุก 4 ชั่วโมง เพื่อนำไปใช้วิเคราะห์ร่วมกับโครงข่ายประสาทเทียม
    \item จงอธิบายกลไกการคำนวณระยะทาง Great-Circle Distance ในคลาส \texttt{GpsWaypoint} และการกรองสัญญาณสะท้อนเทียม (Multipath Noise Filter)
    \item จงอธิบายขั้นตอนการแปลงพิกัดภูมิศาสตร์สู่พิกัดระนาบ 2 มิติของหน้าจอด้วยคลาส \texttt{CustomPainter} ในแอปพลิเคชัน Flutter
\end{enumerate}
\clearpage
